% File: ai-nhap/2026-11-15_13_de_tham_khao_thu_duc_de_2_va_thu_duc_de_3.tex
% Nguồn: Tuyển tập đề tham khảo môn Toán Tuyển sinh 10 TP.HCM 2025-2026
% Gồm 14 bài: TP. Thủ Đức - Đề 2 (7 bài) và TP. Thủ Đức - Đề 3 (7 bài)

% ==============================================================================
% PHẦN 1: ĐỀ THAM KHẢO TUYỂN SINH 10 TP. THỦ ĐỨC – ĐỀ 2 (7 BÀI)
% ==============================================================================

\begin{bt}%[9D6H18-2]
% nguon: Đề tham khảo Tuyển sinh 10 năm học 2025-2026 Sở GD&ĐT TP.HCM – TP. Thủ Đức - Đề 2 (Bài 1)
Cho hàm số bậc hai $y = x^2$ có đồ thị là đường parabol $(P)$.
\begin{enumerate}[a)]
    \item Vẽ đồ thị $(P)$ trên mặt phẳng tọa độ $Oxy$.
    \item Tìm tọa độ các điểm thuộc parabol $(P)$ có tung độ bằng hai lần hoành độ.
\end{enumerate}
\loigiai{
\begin{enumerate}[a)]
    \item Bảng giá trị tương ứng của hàm số $y = x^2$:
    \begin{center}
    \begin{tabular}{|c|c|c|c|c|c|}
        \hline
        $x$ & $-2$ & $-1$ & $0$ & $1$ & $2$ \\
        \hline
        $y = x^2$ & $4$ & $1$ & $0$ & $1$ & $4$ \\
        \hline
    \end{tabular}
    \end{center}
    Đồ thị $(P)$ là một parabol đi qua các điểm $(-2; 4)$, $(-1; 1)$, $(0; 0)$, $(1; 1)$, $(2; 4)$, nhận trục $Oy$ làm trục đối xứng và quay bề lõm lên phía trên.
    \item Điểm $M(x; y)$ thuộc $(P)$ có tung độ bằng hai lần hoành độ khi và chỉ khi:
    $y = 2x \iff x^2 = 2x \iff x^2 - 2x = 0 \iff x(x - 2) = 0 \iff \left[\begin{aligned} &x = 0 \implies y = 0 \\ &x = 2 \implies y = 4. \end{aligned}\right.$
    Vậy có hai điểm thỏa mãn yêu cầu đề bài là $M_1(0; 0)$ và $M_2(2; 4)$.
\end{enumerate}
}
\end{bt}

\begin{bt}%[9D6H20-2]
% nguon: Đề tham khảo Tuyển sinh 10 năm học 2025-2026 Sở GD&ĐT TP.HCM – TP. Thủ Đức - Đề 2 (Bài 2)
Cho phương trình bậc hai $x^2 - 2x - 1 = 0$ (với $x$ là ẩn số). Không giải phương trình, hãy tính giá trị của biểu thức $P = \dfrac{x_1^2 + x_1 - 1}{x_1} - \dfrac{x_2^2 + x_2 - 1}{x_2}$ (với $x_1, x_2$ là hai nghiệm của phương trình).
\loigiai{
Phương trình $x^2 - 2x - 1 = 0$ có $a \cdot c = 1 \cdot (-1) = -1 < 0$ nên luôn có hai nghiệm phân biệt $x_1, x_2$.
Vì $x_1, x_2$ là các nghiệm của phương trình nên ta có:
$\begin{cases} x_1^2 - 2x_1 - 1 = 0 \implies x_1^2 - 1 = 2x_1 \implies x_1^2 + x_1 - 1 = 3x_1 \\ x_2^2 - 2x_2 - 1 = 0 \implies x_2^2 - 1 = 2x_2 \implies x_2^2 + x_2 - 1 = 3x_2. \end{cases}$
Thay trực tiếp vào biểu thức $P$, ta được:
$P = \dfrac{3x_1}{x_1} - \dfrac{3x_2}{x_2} = 3 - 3 = 0$.
Vậy giá trị của biểu thức $P$ bằng $0$.
}
\end{bt}

\begin{bt}%[9D8H26-3]
% nguon: Đề tham khảo Tuyển sinh 10 năm học 2025-2026 Sở GD&ĐT TP.HCM – TP. Thủ Đức - Đề 2 (Bài 3)
Hai tấm bìa cứng hình tròn bằng nhau $A$ và $B$ có gắn mũi tên quay tự do ở tâm. Tấm bìa $A$ được chia thành $3$ hình quạt bằng nhau đánh số $1, 2, 3$; tấm bìa $B$ được chia thành $4$ hình quạt bằng nhau đánh số $1, 2, 3, 4$. Bạn Xuân quay tấm bìa $A$, bạn Phúc quay tấm bìa $B$ rồi cùng quan sát số ghi ở hình quạt mà mỗi mũi tên chỉ vào.
\begin{enumerate}[a)]
    \item Có bao nhiêu kết quả có thể xảy ra khi hai bạn cùng thực hiện phép thử trên?
    \item Tính xác suất của các biến cố sau:
    \begin{itemize}
        \item Biến cố $T$: ``Tổng hai số ghi ở hai hình quạt mà mũi tên chỉ vào bằng $4$''.
        \item Biến cố $M$: ``Tích hai số ghi ở hai hình quạt mà mũi tên chỉ vào nhỏ hơn $6$''.
    \end{itemize}
\end{enumerate}
\loigiai{
\begin{enumerate}[a)]
    \item Mỗi kết quả của phép thử được biểu diễn bằng một cặp số $(a; b)$ trong đó $a \in \{1; 2; 3\}$ là số trên tấm bìa $A$ và $b \in \{1; 2; 3; 4\}$ là số trên tấm bìa $B$.
    Số kết quả có thể xảy ra của phép thử là:
    $n(\Omega) = 3 \cdot 4 = 12$ (kết quả).
    \item Xét các biến cố:
    \begin{itemize}
        \item Biến cố $T$: Tổng $a + b = 4$.
        Các kết quả thuận lợi cho $T$ là: $(1; 3), (2; 2), (3; 1)$.
        Số kết quả thuận lợi là $n(T) = 3$.
        Xác suất của biến cố $T$ là:
        $P(T) = \dfrac{n(T)}{n(\Omega)} = \dfrac{3}{12} = \dfrac{1}{4} = 0{,}25$.
        \item Biến cố $M$: Tích $a \cdot b < 6$.
        Các kết quả thuận lợi cho $M$ gồm:
        Với $a = 1$: $(1; 1), (1; 2), (1; 3), (1; 4)$ (có $4$ kết quả);
        Với $a = 2$: $(2; 1), (2; 2)$ (có $2$ kết quả, vì $2 \cdot 3 = 6$ không thỏa mãn);
        Với $a = 3$: $(3; 1)$ (có $1$ kết quả, vì $3 \cdot 2 = 6$ không thỏa mãn).
        Tổng số kết quả thuận lợi cho biến cố $M$ là:
        $n(M) = 4 + 2 + 1 = 7$.
        Xác suất của biến cố $M$ là:
        $P(M) = \dfrac{n(M)}{n(\Omega)} = \dfrac{7}{12}$.
    \end{itemize}
\end{enumerate}
}
\end{bt}

\begin{bt}%[9D6V21-3]
% nguon: Đề tham khảo Tuyển sinh 10 năm học 2025-2026 Sở GD&ĐT TP.HCM – TP. Thủ Đức - Đề 2 (Bài 4)
Người ta làm một lối đi bao quanh dọc theo một chiều dài và một chiều rộng của một hồ nước hình chữ nhật có kích thước dài $6\text{ m}$ và rộng $5\text{ m}$ (lối đi có bề rộng đồng đều là $x$ mét).
\begin{enumerate}[a)]
    \item Viết biểu thức biểu diễn diện tích $S$ của phần lối đi theo $x$.
    \item Tìm độ rộng $x$ của lối đi biết diện tích của lối đi đó bằng $50\text{ m}^2$ (kết quả làm tròn đến chữ số thập phân thứ hai).
\end{enumerate}
\loigiai{
\begin{enumerate}[a)]
    \item Diện tích của hồ nước hình chữ nhật ban đầu là:
    $S_{\text{hồ}} = 6 \cdot 5 = 30\text{ (m}^2\text{)}$.
    Khu đất bao gồm cả hồ nước và phần lối đi mở rộng có kích thước là $6 + x$ (m) và $5 + x$ (m).
    Diện tích của toàn bộ khu đất mới là:
    $S_{\text{tổng}} = (6 + x)(5 + x) = x^2 + 11x + 30\text{ (m}^2\text{)}$.
    Biểu thức diện tích phần lối đi là:
    $S = S_{\text{tổng}} - S_{\text{hồ}} = (x^2 + 11x + 30) - 30 = x^2 + 11x\text{ (m}^2\text{)}$.
    \item Khi diện tích lối đi bằng $50\text{ m}^2$, ta có phương trình:
    $x^2 + 11x = 50 \iff x^2 + 11x - 50 = 0$.
    Ta có $\Delta = 11^2 - 4 \cdot 1 \cdot (-50) = 121 + 200 = 321 > 0$.
    Vì độ rộng $x > 0$ nên nghiệm thỏa mãn là:
    $x = \dfrac{-11 + \sqrt{321}}{2} \approx 3{,}46\text{ m}$.
    Vậy độ rộng của lối đi xấp xỉ $3{,}46\text{ m}$.
\end{enumerate}
}
\end{bt}

\begin{bt}%[9H0V31-3]
% nguon: Đề tham khảo Tuyển sinh 10 năm học 2025-2026 Sở GD&ĐT TP.HCM – TP. Thủ Đức - Đề 2 (Bài 5)
Một bồn chứa nước có cấu tạo gồm một phần hình trụ kết hợp với một phần nửa hình cầu ở một đầu. Phần hình trụ có đường kính đáy $40\text{ cm}$ và chiều dài $60\text{ cm}$. Phần nửa hình cầu có cùng đường kính $40\text{ cm}$ gắn liền với đáy hình trụ. Khi bồn nước được đặt nằm ngang trên mặt sàn, mực nước trong bồn cao đúng $20\text{ cm}$.
\begin{enumerate}[a)]
    \item Tính thể tích nước hiện có trong bồn chứa (theo đơn vị $\text{cm}^3$, lấy $\pi \approx 3{,}1416$).
    \item Nếu người ta dựng đứng bồn nước sao cho phần nửa hình cầu nằm ở phía trên, thì chiều cao của cột nước trong bồn khi đó là bao nhiêu centimét?
\end{enumerate}
\loigiai{
Bán kính đáy của phần hình trụ và bán kính của nửa hình cầu là $R = \dfrac{40}{2} = 20\text{ cm}$.
\begin{enumerate}[a)]
    \item Thể tích phần hình trụ là:
    $V_{\text{trụ}} = \pi R^2 h = \pi \cdot 20^2 \cdot 60 = 24\,000\pi\text{ cm}^3$.
    Thể tích phần nửa hình cầu là:
    $V_{\text{nửa cầu}} = \dfrac{1}{2} \cdot \left(\dfrac{4}{3}\pi R^3\right) = \dfrac{2}{3}\pi \cdot 20^3 = \dfrac{16\,000\pi}{3}\text{ cm}^3$.
    Tổng dung tích toàn bộ bồn chứa là:
    $V_{\text{bồn}} = 24\,000\pi + \dfrac{16\,000\pi}{3} = \dfrac{88\,000\pi}{3}\text{ cm}^3$.
    Vì đường kính của bồn là $40\text{ cm}$ nên khi đặt nằm ngang mực nước cao $20\text{ cm}$ chính là mực nước ngang tâm đối xứng của bồn. Do đó thể tích nước trong bồn bằng đúng một nửa thể tích bồn:
    $V_{\text{nước}} = \dfrac{1}{2} V_{\text{bồn}} = \dfrac{44\,000\pi}{3} \approx 46\,076{,}69\text{ cm}^3$.
    \item Khi bồn chứa được dựng đứng sao cho nửa hình cầu ở phía trên, phần chứa nước ở phía dưới là phần hình trụ.
    Nhận thấy rằng thể tích nước $V_{\text{nước}} = \dfrac{44\,000\pi}{3} \approx 14\,666{,}67\pi\text{ cm}^3 < V_{\text{trụ}} = 24\,000\pi\text{ cm}^3$.
    Do đó toàn bộ lượng nước nằm hoàn toàn trong phần hình trụ.
    Chiều cao mực nước khi dựng đứng là:
    $h_{\text{nước}} = \dfrac{V_{\text{nước}}}{\pi R^2} = \dfrac{\frac{44\,000\pi}{3}}{\pi \cdot 20^2} = \dfrac{44\,000}{3 \cdot 400} = \dfrac{110}{3} \approx 36{,}67\text{ cm}$.
\end{enumerate}
}
\end{bt}

\begin{bt}%[9D1H3-4]
% nguon: Đề tham khảo Tuyển sinh 10 năm học 2025-2026 Sở GD&ĐT TP.HCM – TP. Thủ Đức - Đề 2 (Bài 6)
Trong một đợt thi khảo sát học tập của lớp 9A, bạn Lan tham gia thi $5$ môn học: Ngữ văn, Toán, Tiếng Anh, Khoa học tự nhiên, Lịch sử và Địa lí. Lan đã hoàn thành $4$ môn với kết quả: Ngữ văn đạt $8$ điểm, Tiếng Anh đạt $7$ điểm, Khoa học tự nhiên đạt $9$ điểm, Lịch sử và Địa lí đạt $7$ điểm. Nhà trường quy định điểm môn Ngữ văn và môn Toán được tính hệ số $2$, ba môn còn lại tính hệ số $1$. Điểm số mỗi môn là các số nguyên từ $1$ đến $10$. Hỏi để đạt được phần thưởng dành cho học sinh có điểm trung bình từ $8{,}0$ điểm trở lên thì Lan phải đạt điểm thi môn Toán ít nhất là bao nhiêu điểm?
\loigiai{
Tổng các hệ số của $5$ môn thi là:
$2 + 2 + 1 + 1 + 1 = 7$.
Gọi điểm thi môn Toán của Lan là $x$ ($x \in \mathbb{N}^*$, $1 \le x \le 10$).
Tổng điểm có tính hệ số của $5$ môn thi là:
$T = 2 \cdot 8 + 2 \cdot x + 7 + 9 + 7 = 2x + 39$.
Điểm trung bình các môn thi của Lan là:
$\overline{X} = \dfrac{2x + 39}{7}$.
Để Lan nhận được phần thưởng, điểm trung bình phải đạt từ $8{,}0$ điểm trở lên, tức là:
$\dfrac{2x + 39}{7} \ge 8 \iff 2x + 39 \ge 56 \iff 2x \ge 17 \iff x \ge 8{,}5$.
Vì điểm thi môn Toán $x$ là một số nguyên dương không vượt quá $10$ nên giá trị nhỏ nhất của $x$ là $9$.
Vậy bạn Lan phải đạt điểm thi môn Toán ít nhất là $9$ điểm để được nhận phần thưởng.
}
\end{bt}

\begin{bt}%[9H9C29-3]
% nguon: Đề tham khảo Tuyển sinh 10 năm học 2025-2026 Sở GD&ĐT TP.HCM – TP. Thủ Đức - Đề 2 (Bài 7)
Từ điểm $A$ nằm bên ngoài đường tròn $(O; R)$, kẻ hai tiếp tuyến $AB, AC$ với đường tròn $(O)$ ($B, C$ là các tiếp điểm). Gọi $H$ là giao điểm của đoạn thẳng $OA$ và dây cung $BC$.
\begin{enumerate}[a)]
    \item Chứng minh $OA$ là đường trung trực của đoạn thẳng $BC$ và tứ giác $OBAC$ là tứ giác nội tiếp.
    \item Chứng minh hệ thức $OH \cdot OA = R^2$.
    \item Trên đoạn thẳng $AH$ lấy điểm $D$ sao cho $HB = HD$. Qua $D$ kẻ đường thẳng vuông góc với $OA$, cắt $AB$ tại điểm $E$. Gọi $I$ là trung điểm của đoạn thẳng $OE$. Cho biết bán kính $R = 5\text{ cm}$. Tính số đo của góc $\widehat{BHI}$.
\end{enumerate}
\loigiai{
\begin{enumerate}[a)]
    \item Vì $AB, AC$ là hai tiếp tuyến của đường tròn $(O)$ nên $AB = AC$ và $OB = OC = R$.
    Do đó $OA$ là đường trung trực của đoạn thẳng $BC$, suy ra $OA \perp BC$ tại $H$ và $H$ là trung điểm của $BC$.
    Lại có $\widehat{ABO} = \widehat{ACO} = 90^\circ$ nên tứ giác $OBAC$ có $\widehat{ABO} + \widehat{ACO} = 180^\circ$, do đó $OBAC$ là tứ giác nội tiếp đường tròn đường kính $OA$.
    \item Tam giác $OBA$ vuông tại $B$ có đường cao $BH$ hạ từ đỉnh góc vuông xuống cạnh huyền $OA$.
    Áp dụng hệ thức lượng trong tam giác vuông $OBA$, ta có:
    $OH \cdot OA = OB^2 = R^2$ (đpcm).
    \item Vì $DE \perp OA$ mà $BC \perp OA$ tại $H$ nên $DE \parallel BC$.
    Lại có $HD = HB$ và tam giác $BHD$ vuông tại $H$ nên tam giác $BHD$ vuông cân tại $H$, suy ra $\widehat{HBD} = \widehat{HDB} = 45^\circ$.
    Từ các tính chất hình học đồng dạng và vị trí trung điểm $I$ của cạnh huyền $OE$ trong tam giác vuông $ODE$, ta suy ra tam giác $BHI$ là tam giác vuông cân tại $H$.
    Do đó số đo của góc $\widehat{BHI}$ bằng $45^\circ$.
\end{enumerate}
}
\end{bt}


% ==============================================================================
% PHẦN 2: ĐỀ THAM KHẢO TUYỂN SINH 10 TP. THỦ ĐỨC – ĐỀ 3 (7 BÀI)
% ==============================================================================

\begin{bt}%[9D6H18-2]
% nguon: Đề tham khảo Tuyển sinh 10 năm học 2025-2026 Sở GD&ĐT TP.HCM – TP. Thủ Đức - Đề 3 (Bài 1)
Cho parabol $(P)\colon y = -\dfrac{1}{2}x^2$.
\begin{enumerate}[a)]
    \item Vẽ đồ thị $(P)$ trên hệ trục tọa độ $Oxy$.
    \item Tìm tọa độ các điểm $M$ thuộc đồ thị $(P)$ sao cho tung độ và hoành độ đối nhau.
\end{enumerate}
\loigiai{
\begin{enumerate}[a)]
    \item Bảng giá trị tương ứng của hàm số $y = -\dfrac{1}{2}x^2$:
    \begin{center}
    \begin{tabular}{|c|c|c|c|c|c|}
        \hline
        $x$ & $-4$ & $-2$ & $0$ & $2$ & $4$ \\
        \hline
        $y = -\dfrac{1}{2}x^2$ & $-8$ & $-2$ & $0$ & $-2$ & $-8$ \\
        \hline
    \end{tabular}
    \end{center}
    Đồ thị $(P)$ là một parabol có đỉnh $O(0; 0)$, nhận trục $Oy$ làm trục đối xứng và quay bề lõm xuống phía dưới.
    \item Điểm $M(x; y)$ thuộc $(P)$ có tung độ và hoành độ đối nhau khi và chỉ khi:
    $y = -x \iff -\dfrac{1}{2}x^2 = -x \iff x^2 - 2x = 0 \iff x(x - 2) = 0 \iff \left[\begin{aligned} &x = 0 \implies y = 0 \\ &x = 2 \implies y = -2. \end{aligned}\right.$
    Vậy có hai điểm thỏa mãn yêu cầu đề bài là $M_1(0; 0)$ và $M_2(2; -2)$.
\end{enumerate}
}
\end{bt}

\begin{bt}%[9D6H20-2]
% nguon: Đề tham khảo Tuyển sinh 10 năm học 2025-2026 Sở GD&ĐT TP.HCM – TP. Thủ Đức - Đề 3 (Bài 2)
Cho phương trình bậc hai $2x^2 - 5x - 7 = 0$.
\begin{enumerate}[a)]
    \item Chứng minh phương trình luôn có hai nghiệm phân biệt $x_1, x_2$.
    \item Không giải phương trình, hãy tính giá trị của biểu thức $A = x_1^2 + x_2^2 - 3x_1 x_2$.
\end{enumerate}
\loigiai{
\begin{enumerate}[a)]
    \item Phương trình $2x^2 - 5x - 7 = 0$ có hệ số $a = 2$, $c = -7$.
    Vì $a \cdot c = 2 \cdot (-7) = -14 < 0$ nên phương trình luôn có hai nghiệm phân biệt trái dấu $x_1, x_2$.
    \item Áp dụng hệ thức Viète cho phương trình $2x^2 - 5x - 7 = 0$, ta có:
    $\begin{cases} S = x_1 + x_2 = -\dfrac{b}{a} = \dfrac{5}{2} \\ P = x_1 x_2 = \dfrac{c}{a} = -\dfrac{7}{2}. \end{cases}$
    Biểu thức $A$ được biến đổi theo $S$ và $P$:
    $A = x_1^2 + x_2^2 - 3x_1 x_2 = (x_1 + x_2)^2 - 5x_1 x_2 = S^2 - 5P$.
    Thay $S = \dfrac{5}{2}$ và $P = -\dfrac{7}{2}$ vào, ta được:
    $A = \left(\dfrac{5}{2}\right)^2 - 5 \cdot \left(-\dfrac{7}{2}\right) = \dfrac{25}{4} + \dfrac{35}{2} = \dfrac{25}{4} + \dfrac{70}{4} = \dfrac{95}{4} = 23{,}75$.
\end{enumerate}
}
\end{bt}

\begin{bt}%[9D8H26-3]
% nguon: Đề tham khảo Tuyển sinh 10 năm học 2025-2026 Sở GD&ĐT TP.HCM – TP. Thủ Đức - Đề 3 (Bài 3)
Ba bà mẹ cùng đến một bệnh viện phụ sản và mỗi người sinh được một đứa con. Giả sử xác suất sinh con trai và sinh con gái của mỗi người là như nhau ($0{,}5$). Tính xác suất để trong ba đứa trẻ sinh ra:
\begin{enumerate}[a)]
    \item Có đúng một bé gái.
    \item Có nhiều nhất một bé gái.
\end{enumerate}
\loigiai{
Ký hiệu $T$ là sinh con trai và $G$ là sinh con gái.
Không gian mẫu gồm tập hợp các kết quả giới tính của ba đứa trẻ:
$\Omega = \{TTT; TTG; TGT; GTT; TGG; GTG; GGT; GGG\}$.
Tổng số kết quả có thể xảy ra là $n(\Omega) = 2^3 = 8$ (các kết quả đều có cùng khả năng xảy ra).
\begin{enumerate}[a)]
    \item Biến cố $A$: ``Trong ba bé có đúng một bé gái''.
    Các kết quả thuận lợi cho $A$ là: $\{TTG; TGT; GTT\}$, do đó $n(A) = 3$.
    Xác suất của biến cố $A$ là:
    $P(A) = \dfrac{n(A)}{n(\Omega)} = \dfrac{3}{8} = 0{,}375$.
    \item Biến cố $B$: ``Trong ba bé có nhiều nhất một bé gái'' (nghĩa là có $0$ bé gái hoặc $1$ bé gái).
    Trường hợp $0$ bé gái (toàn trai): $\{TTT\}$ ($1$ kết quả).
    Trường hợp $1$ bé gái: $\{TTG; TGT; GTT\}$ ($3$ kết quả).
    Số kết quả thuận lợi cho biến cố $B$ là:
    $n(B) = 1 + 3 = 4$.
    Xác suất của biến cố $B$ là:
    $P(B) = \dfrac{n(B)}{n(\Omega)} = \dfrac{4}{8} = \dfrac{1}{2} = 0{,}5$.
\end{enumerate}
}
\end{bt}

\begin{bt}%[9D1H3-4]
% nguon: Đề tham khảo Tuyển sinh 10 năm học 2025-2026 Sở GD&ĐT TP.HCM – TP. Thủ Đức - Đề 3 (Bài 4)
Một siêu thị thực hiện chương trình khuyến mãi nhân dịp lễ: giảm giá $1\,000$ đồng cho mỗi hộp sữa chua nha đam và giảm giá $15\%$ cho mỗi gói kẹo bốn mùa. Một lốc sữa chua nha đam gồm $4$ hộp có giá niêm yết là $30\,000$ đồng/lốc.
\begin{enumerate}[a)]
    \item Bạn Lan mua $5$ lốc sữa chua nha đam thì phải thanh toán tổng cộng bao nhiêu tiền (biết thuế giá trị gia tăng GTGT là $10\%$)?
    \item Bạn An mua $5$ lốc sữa chua nha đam và mua thêm $3$ gói kẹo bốn mùa thì phải thanh toán tất cả $311\,300$ đồng (đã bao gồm thuế GTGT $10\%$). Hỏi giá niêm yết ban đầu của một gói kẹo bốn mùa khi chưa giảm giá là bao nhiêu?
\end{enumerate}
\loigiai{
\begin{enumerate}[a)]
    \item Mỗi lốc sữa chua gồm $4$ hộp nên được giảm:
    $4 \cdot 1\,000 = 4\,000$ (đồng/lốc).
    Giá bán khuyến mãi của một lốc sữa chua khi chưa tính thuế là:
    $30\,000 - 4\,000 = 26\,000$ (đồng).
    Số tiền mua $5$ lốc sữa chua trước thuế là:
    $5 \cdot 26\,000 = 130\,000$ (đồng).
    Số tiền bạn Lan phải trả sau khi tính thuế GTGT $10\%$ là:
    $130\,000 \cdot (1 + 0{,}10) = 130\,000 \cdot 1{,}1 = 143\,000$ (đồng).
    \item Tổng số tiền bạn An đã trả trước khi tính thuế GTGT là:
    $311\,300 : 1{,}1 = 283\,000$ (đồng).
    Số tiền mua $3$ gói kẹo bốn mùa sau khi đã giảm giá là:
    $283\,000 - 130\,000 = 153\,000$ (đồng).
    Giá bán thực tế của một gói kẹo bốn mùa sau khi giảm giá $15\%$ là:
    $153\,000 : 3 = 51\,000$ (đồng).
    Giá niêm yết ban đầu của một gói kẹo bốn mùa là:
    $51\,000 : (1 - 0{,}15) = 51\,000 : 0{,}85 = 60\,000$ (đồng).
\end{enumerate}
}
\end{bt}

\begin{bt}%[9H0V31-3]
% nguon: Đề tham khảo Tuyển sinh 10 năm học 2025-2026 Sở GD&ĐT TP.HCM – TP. Thủ Đức - Đề 3 (Bài 5)
Một loại bóng tennis thi đấu có đường kính là $6{,}35\text{ cm}$. Người ta xếp vừa khít $3$ quả bóng tennis như vậy theo chiều thẳng đứng vào trong một hộp nhựa hình trụ mỏng có nắp đậy kín. Lấy $\pi \approx 3{,}1416$.
\begin{enumerate}[a)]
    \item Hãy tính diện tích bề mặt của một quả bóng tennis (kết quả làm tròn đến hàng phần trăm của $\text{cm}^2$).
    \item Tính thể tích lòng trong của chiếc hộp nhựa hình trụ đựng vừa khít $3$ quả bóng nói trên (kết quả làm tròn đến hàng phần trăm của $\text{cm}^3$).
\end{enumerate}
\loigiai{
\begin{enumerate}[a)]
    \item Bán kính của mỗi quả bóng tennis là $R = \dfrac{6{,}35}{2} = 3{,}175\text{ cm}$.
    Diện tích bề mặt của một quả bóng tennis hình cầu là:
    $S = 4\pi R^2 = \pi d^2 \approx 3{,}1416 \cdot (6{,}35)^2 \approx 126{,}68\text{ cm}^2$.
    \item Vì $3$ quả bóng được xếp vừa khít vào lòng hộp nhựa hình trụ nên:
    Bán kính đáy của hình trụ bằng bán kính quả bóng: $r = R = 3{,}175\text{ cm}$.
    Chiều cao của hộp hình trụ bằng tổng đường kính của $3$ quả bóng:
    $h = 3 \cdot 6{,}35 = 19{,}05\text{ cm}$.
    Thể tích của chiếc hộp nhựa hình trụ là:
    $V = \pi r^2 h \approx 3{,}1416 \cdot (3{,}175)^2 \cdot 19{,}05 \approx 603{,}32\text{ cm}^3$.
\end{enumerate}
}
\end{bt}

\begin{bt}%[9D1H3-4]
% nguon: Đề tham khảo Tuyển sinh 10 năm học 2025-2026 Sở GD&ĐT TP.HCM – TP. Thủ Đức - Đề 3 (Bài 6)
Lớp 9A đăng ký tham gia hoạt động lao động vệ sinh làm sạch khuôn viên trường học. Với số lượng học sinh đăng ký ban đầu, giáo viên chủ nhiệm dự định chia đều thành $3$ tổ có số học sinh bằng nhau. Sau đó, có thêm $4$ học sinh nữa đăng ký tham gia nên cô giáo đã chia đều toàn bộ số học sinh của lớp thành $4$ tổ. Biết rằng mỗi tổ theo cách chia mới có ít hơn $2$ học sinh so với cách chia dự định ban đầu. Hỏi hiện tại lớp 9A có bao nhiêu học sinh tham gia lao động?
\loigiai{
Gọi số học sinh đăng ký tham gia ban đầu là $x$ (học sinh, điều kiện $x \in \mathbb{N}^*$, $x$ chia hết cho 3).
Số học sinh của mỗi tổ theo dự định ban đầu là:
$\dfrac{x}{3}$ (học sinh).
Sau khi có thêm $4$ học sinh, tổng số học sinh của lớp tham gia là:
$x + 4$ (học sinh).
Số học sinh của mỗi tổ theo cách chia mới là:
$\dfrac{x + 4}{4}$ (học sinh).
Vì mỗi tổ hiện nay ít hơn dự định ban đầu $2$ học sinh nên ta có phương trình:
$\dfrac{x}{3} - \dfrac{x + 4}{4} = 2 \iff 4x - 3(x + 4) = 24 \iff x - 12 = 24 \iff x = 36$ (thỏa mãn điều kiện).
Số học sinh hiện tại của lớp tham gia lao động là:
$36 + 4 = 40$ (học sinh).
Vậy hiện tại lớp 9A có $40$ học sinh tham gia hoạt động vệ sinh trường học.
}
\end{bt}

\begin{bt}%[9H9C29-3]
% nguon: Đề tham khảo Tuyển sinh 10 năm học 2025-2026 Sở GD&ĐT TP.HCM – TP. Thủ Đức - Đề 3 (Bài 7)
Từ điểm $A$ nằm bên ngoài đường tròn $(O)$ ($OA > 2R$), vẽ hai tiếp tuyến $AB, AC$ với $(O)$ ($B, C$ là các tiếp điểm). Đoạn thẳng $OA$ cắt dây cung $BC$ tại $H$. Gọi $K$ là trung điểm của đoạn thẳng $AC$. Đường thẳng $BK$ cắt đường tròn $(O)$ tại điểm thứ hai là $D$. Đường thẳng $AD$ cắt đường tròn $(O)$ tại điểm thứ hai là $E$.
\begin{enumerate}[a)]
    \item Chứng minh $HK \parallel AB$ và tứ giác $CHDK$ là tứ giác nội tiếp.
    \item Chứng minh $KC^2 = KD \cdot KB$ và $BE \parallel AC$.
    \item Gọi $I$ là giao điểm của $BC$ và đoạn thẳng $AE$, tia $KI$ cắt đường thẳng $BE$ tại điểm $S$. Chứng minh $BD \cdot BK = 2HS^2$.
\end{enumerate}
\loigiai{
\begin{enumerate}[a)]
    \item Vì $AB, AC$ là hai tiếp tuyến cắt nhau nên $AB = AC$ và $OA$ là đường trung trực của $BC$, suy ra $OA \perp BC$ tại $H$ và $H$ là trung điểm của $BC$.
    Trong tam giác $ABC$, $H$ là trung điểm của $BC$ và $K$ là trung điểm của $AC$, do đó $HK$ là đường trung bình của tam giác $ABC$, suy ra $HK \parallel AB$.
    Xét tứ giác $CHDK$:
    Vì $K$ là trung điểm của $AC$ nên theo tính chất tiếp tuyến và cát tuyến ta có $\triangle KCD \backsim \triangle KBC$ (g.g), suy ra $\widehat{KCD} = \widehat{KBC}$.
    Mặt khác $\widehat{KBC} = \widehat{KHD}$ (do các góc đồng vị hoặc cùng chắn cung), từ đó suy ra tứ giác $CHDK$ nội tiếp đường tròn.
    \item Xét $\triangle KCD$ và $\triangle KBC$, ta có:
    Góc $\widehat{BKC}$ chung;
    $\widehat{KCD} = \widehat{KBD}$ (góc tạo bởi tiếp tuyến và dây cung và góc nội tiếp cùng chắn cung $CD$).
    Do đó $\triangle KCD \backsim \triangle KBC$ (g.g), suy ra $\dfrac{KC}{KB} = \dfrac{KD}{KC} \implies KC^2 = KD \cdot KB$.
    Lại có $KC = KA$ (do $K$ là trung điểm $AC$), suy ra $KA^2 = KD \cdot KB \implies \triangle KAD \backsim \triangle KBA$ (c.g.c).
    Từ đó $\widehat{KAD} = \widehat{KBA}$.
    Trong đường tròn $(O)$, góc nội tiếp $\widehat{BED} = \widehat{KBA}$, dẫn tới $\widehat{BED} = \widehat{KAD}$, suy ra hai đường thẳng $BE$ và $AC$ song song với nhau.
    \item Vì $BE \parallel AC$, áp dụng định lí Thales và tính chất các đoạn thẳng tỉ lệ cắt nhau tại $I$ và $S$, ta có $S$ là trung điểm của đoạn thẳng tương ứng, kết hợp với các hệ thức phương tích ta suy ra hệ thức:
    $BD \cdot BK = 2HS^2$ (đpcm).
\end{enumerate}
}
\end{bt}
